% answers_short.tex —— 2000 年·上海卷·理科 · 简短答案速查
% 由 tools/make_short_answers.py 自 answers.tex 的 \gksolhead/\gkanswer 提取，
% 并与 2000/2000_数学真题_上海理_答案要点.md 逐题比对后生成，勿手改（改源 MD 后重跑）。

\begin{gknotice}
\textbf{使用说明}：本节只给最终答案便于快速核对，完整推导见解析版。答案为编写者独立推导，\textbf{非官方标准答案}。
\end{gknotice}

\gkpart{填空题}
\begin{itemize}[leftmargin=2.6em]
  \item[\textbf{1.}] $m=\dfrac32$。
  \item[\textbf{2.}] $\left(\dfrac12,3\right)$。
  \item[\textbf{3.}] 焦点为 $(-4,0)$、$(6,0)$。
  \item[\textbf{4.}] $e^{-2}$。
  \item[\textbf{5.}] $b=1$。
  \item[\textbf{6.}] 至少 $9$ 年。
  \item[\textbf{7.}] 可填“三条侧棱相等”。
  \item[\textbf{8.}] $f(x)=x$（$1\leq x\leq2$）。
  \item[\textbf{9.}] $-462$。
  \item[\textbf{10.}] $\dfrac1{14}$。
  \item[\textbf{11.}] $2\sqrt3$。
  \item[\textbf{12.}] U/PENDING。
\end{itemize}

\gkpart{选择题}
\begin{center}\begin{tabular}{ccccc}
\toprule
题号 & 13 & 14 & 15 & 16 \\
\midrule
答案 & C & A & A & D \\
\bottomrule
\end{tabular}\end{center}

\gkpart{解答题}
\begin{itemize}[leftmargin=2.6em]
  \item[\textbf{17.}] 中点坐标为 $\left(-\dfrac95,\dfrac15\right)$。
  \item[\textbf{18.}] 四面体体积为 $\dfrac83$。
  \item[\textbf{19.}] (1) $\dfrac72$，在 $x=1$ 处取得；(2) $a>-3$
  \item[\textbf{20.}] (1) 指令为 $(4\sqrt2,45^\circ)$；(2) 最快在 $(7,0)$ 截住小球，所需指令为 $(5.00,-98.13^\circ)$
  \item[\textbf{21.}] (1) 按原页指数 $2$ 的字面推导，$b_n=20a^2$；(2) $0<a<10$；(3) U/PENDING
  \item[\textbf{22.}] (1) $m=\sqrt3$，$x'=x+\sqrt3y$，$y'=\sqrt3x-y$；(2) $y'=(2-\sqrt3)x'+2-2\sqrt3$；(3) $y=\dfrac{x}{\sqrt3}$ 或 $y=-\sqrt3x$
\end{itemize}
